\section{引\quad 言}

纳米材料的宏观性能并不只取决于化学组成，而在很大程度上由晶型、晶粒尺寸、形貌与表面状态共同决定；这些结构参数又可以在制备与后处理阶段被系统调控。在众多调控手段中，热处理最为简便：它既能提供相变所需的能量，又会驱动晶粒长大与颗粒烧结，因而是连接“制备”与“性能”的关键环节\cite{wang2005表征,li2008制备}。

二氧化钛是研究这一问题的理想模型体系。它存在锐钛矿、金红石和板钛矿三种晶型，其中锐钛矿为间接带隙半导体（约 3.2 eV），金红石为直接带隙半导体（约 3.0 eV）\cite{chen2007tio2}。热力学上金红石更稳定，但纳米尺度下表面能与晶粒尺寸会改变相稳定性，使相变温度随前驱体、杂质含量、气氛与升温制度在较宽区间内移动\cite{zhang2000phase}。同时，混相（锐钛矿/金红石）界面有利于电荷分离，往往比单一相具有更高的光催化活性\cite{ohtani1997mixture}，因此定量描述相变进程与晶粒长大规律具有实际意义。

从制备角度看，溶胶-凝胶法通过金属醇盐水解-缩聚获得均匀的前驱体，产物纯度高、组分易调，适合制备纳米颗粒；静电纺丝法则借助高压电场将溶胶拉伸成连续纤维，可获得高长径比、易回收的一维材料\cite{zhao2019静电纺丝}。两种路线得到的都是 TiO$_2$ 前驱体（水合氧化物/含聚合物复合纤维），后续热处理条件基本相同，这为比较“成型方式”与“热处理历史”对最终结构的影响提供了便利。

本文以第一轮专业实验中获得的 TiO$_2$ 系列样品为对象，选取溶胶-凝胶制备、静电纺丝制备、X 射线衍射、扫描电子显微镜与能谱、拉曼光谱、紫外-可见漫反射六个实验项目的数据开展分析。工作重点不在于复述各方法的原理，而是用同一温度序列（200—900 $^\circ$C）把四类表征结果串联起来，回答三个问题：(1) 相变发生在什么温度区间，转变程度如何随温度演化；(2) 晶粒与颗粒尺寸如何随温度长大，两者差异反映了什么结构特征；(3) 结构与形貌的变化如何体现在光学带隙上。通过交叉验证，期望给出一个可自洽解释的实验图像，并讨论其中的误差来源。
